\ExerciseSection

\begin{enumerate}
\def\labelenumi{\arabic{enumi})}
\item
  Let \(a\) be a non-zero complex number, and \(n\) an integer \(>0\) . For each real number \(r > 0\) such that \(r^n \leqslant |a|\) , show that there exists \(z \in \mathbf{C}\) such that \(|z| = r\) and \(|a + z^n| = |a| - r^n\) . Deduce that, if \(f\) is a polynomial of degree \(>0\) , with complex coefficients, we cannot have \(|f(z_0)| \leqslant |f(z)|\) at all points \(z\) of a neighbourhood of a point \(z_0\) , provided that \(f(z_0) \neq 0\) .
\item
  Show that, if \(f\) is any polynomial with complex coefficients, not identically zero, there exists a real number \(r > 0\) such that
\end{enumerate}

\[
\left| f (z) \right| > \left| f (\mathbf {o}) \right|
\]

whenever \(|z| \geqslant r\) . Hence, with the help of Exercise 1 and Weierstrass' theorem (Chapter IV, § 6, no. 1, Theorem 1), give another proof of the fact that \(\mathbf{C}\) is algebraically closed (consider the function \(f\) on the compact set of points \(z\) such that \(|z| \leqslant r\) ).

\begin{enumerate}
\def\labelenumi{\arabic{enumi})}
\setcounter{enumi}{2}
\item
  Show, without using Theorem 1, that the mapping \(z \to z / \overline{z}\) is a strict morphism of the topological group \(\mathbf{C}^*\) onto the topological group \(\mathbf{U}\) . Deduce that the mapping \(t \to (\mathrm{i} + it) / (\mathrm{i} - it) \left[ (\mathrm{i} + it) / (\mathrm{i} - it) = -\mathrm{i} \right.\) if \(t = \infty\) is an isomorphism of the topological group \(\tilde{\mathbf{R}}\) (no. 6) onto the topological group \(\mathbf{U}\) , and hence give another proof of Theorem 1.
\item
  Let K be the smallest Pythagorean subfield of R and let \(\mathrm{K}^{\prime}=\mathrm{K}(i)\) . Let G be the multiplicative group of elements of \(K^{\prime}\) of absolute value i (a subgroup of U). Show that G is not isomorphic to the additive group of numbers of K mod i. (Observe that in the latter group there exist elements of any prime order p; on the other hand, if p is a prime such that p-1 is not a power of 2, show that G contains no pth root of unity other than 1, by noting that the degree over Q of every element of \(K^{\prime}\) is a power of 2.)
\end{enumerate}

